% !TEX root = ../main.tex



\section{Experiments}
\label{sec:experiments}

\subsection{Experimental Setup}
\label{sec:setup}

\noindent\textbf{Benchmarks.} We evaluate on three open-ended task streams (Table~\ref{tab:bench_stats}): \emph{PolyBench} for prediction markets~\cite{cheng2026polybench}, \emph{CTF-Dojo} for security challenges~\cite{zhuo2025training}, and \emph{FutureX} for event forecasting~\cite{zeng2025futurex}. All three enforce strict temporal order and covers the three dimensions of challenge. Details and non-stationarity diagnostics are in Appendix~\ref{app:benchmarks} and~\ref{app:nonstationarity}.

\begin{table}[!t]
\centering
\footnotesize
\setlength{\tabcolsep}{2pt}
\caption{Benchmark statistics for the three chronological task streams.}
\vskip -1em
\label{tab:bench_stats}
\resizebox{\columnwidth}{!}{%
\begin{tabular}{@{}lrlll@{}}
\toprule
Bench. & Tasks & Span &Domain \\
\midrule
PolyBench & 5{,}075 & Feb 6--22, 2026 & Prediction markets \\
CTF-Dojo  & 261   & 2011--2024  & Security \\
FutureX   & 503   & Jan--Apr 2026 & Forecasting \\
\bottomrule
\end{tabular}
}
\vskip -1em
\end{table}


\begin{table*}[!t]
\centering
\caption{Main comparison across no-evolution agents, auto-harness baselines, a human-designed system, and our three variants. PolyBench reports Accuracy / Return (coverage-scaled CWR, in \%); CTF-Dojo and FutureX report the official Pass@1. Bold marks the best result per row and underline marks second best.}
\vskip -0.5em
\label{tab:rq1_main}
\begingroup
\small
\setlength{\tabcolsep}{3pt}
\def\arraystretch{1.12}
\resizebox{\textwidth}{!}{%
\begin{tabular}{@{}llccccccccccccccc@{}}
\toprule
\multicolumn{2}{c}{} &
\multicolumn{5}{c}{\textbf{No evolution}} &
\multicolumn{5}{c}{\textbf{Auto-harness baselines}} &
\multicolumn{1}{c}{\textbf{Human}} &
\multicolumn{3}{c}{\textbf{Ours}} \\
\cmidrule(lr){3-7} \cmidrule(lr){8-12} \cmidrule(lr){13-13} \cmidrule(lr){14-16}
\textbf{Benchmark} & \textbf{Metric} &
\textbf{Sonnet} &
\textbf{Haiku} &
\textbf{DeepSeek} &
\textbf{Kimi} &
\textbf{GLM} &
\textbf{A-Evolve} &
\textbf{GEPA} &
\makecell{\textbf{Meta}\\\textbf{Harness}} &
\makecell{\textbf{Cont.}\\\textbf{Harness}} &
\textbf{SkillOS} &
\textbf{OctoTools} &
\makecell{\textbf{Multi}\\\textbf{agent}} &
\textbf{Adaptive} &
\makecell{\textbf{Full}\\\textbf{System}} \\
\midrule
\multirow{2}{*}{\textbf{PolyBench}}
& Accuracy {\scriptsize$\uparrow$} & 22.2 & 15.2 & 1.4 & 14.0 & 14.0 & 18.4 & 13.4 & 50.8 & 8.5 & 21.4 & 40.0 & \underline{79.8} & 77.4 & \textbf{80.9} \\
& Return {\scriptsize$\uparrow$} & $+$1.7 & $+$88.0 & $+$16.4 & $-$2.0 & $+$10.1 & $+$7.2 & $+$0.2 & $+$320 & $+$1.7 & $+$3.6 & $+$20.4 & \underline{$+$351} & \textbf{$+$352} & $+$330 \\
\midrule
\textbf{CTF-Dojo} & Pass {\scriptsize$\uparrow$} & 37.2 & 23.8 & 26.1 & 24.5 & 12.6 & 45.2 & 42.9 & 41.0 & 25.7 & 29.5 & 38.3 & \underline{47.9} & 46.0 & \textbf{50.2} \\
\midrule
\textbf{FutureX} & Pass {\scriptsize$\uparrow$} & 31.0 & 31.0 & 31.2 & 27.8 & 30.8 & \underline{47.5} & 28.2 & 29.4 & 31.8 & 29.8 & 25.6 & \textbf{49.5} & 44.1 & 47.3 \\
\bottomrule
\end{tabular}
}
\endgroup
\vskip -0.7em
\end{table*}

\noindent\textbf{Baselines.} We compare with no-evolution runs with different solver agents using \textbf{Sonnet-4.6}~\cite{anthropic2025claude46}, \textbf{DeepSeek-V3.2}~\cite{deepseek2024v3}, \textbf{Claude Haiku-4.5}~\cite{anthropic2025haiku45}, \textbf{GLM-4.7}~\cite{glm2024chatglm}, and \textbf{Kimi-K2.5}~\cite{moonshot2025kimik2}. We compare with five auto-harness baselines: \textbf{A-Evolve}~\cite{lin2026position}, \textbf{GEPA}~\cite{agrawal2025gepa}, \textbf{Meta-Harness}~\cite{lee2026meta}, \textbf{Continual Harness}~\cite{karten2026continual}, and \textbf{SkillOS}~\cite{ouyang2026skillos}. We also compare with one human-designed system \textbf{OctoTools}~\cite{lu2025octotools}.


\noindent\textbf{Solver and evolver.} We use Claude Sonnet 4.6 as the solver for all experiments unless specified, and use Claude Opus 4.6 as the evolver, both at temperature 0 to attribute gains to the evolution algorithm rather than sampling noise; the no-evolution controls additionally report base-agent results for Haiku 4.5, DeepSeek-V3.2, Kimi-K2.5, and GLM-4.7. All algorithms share the same batch size (100/20/20 for PolyBench/CTF-Dojo/FutureX), batch loop, and temporal-reveal gate; only the evolution algorithm varies.

\noindent\textbf{Metrics.} For PolyBench, we report two complementary metrics: \emph{Accuracy}, the fraction of all markets traded correctly; and \emph{Return}, defined as $\mathrm{Coverage}\times\mathrm{CWR}$, where CWR is the confidence-weighted portfolio profit-to-investment ratio over traded markets and coverage is the fraction of traded. For CTF-Dojo and FutureX, we report the official pass rate (Pass@1) defined by the original benchmarks. We also report \emph{lift} in the figures as the difference with baselines. Full definitions are in Appendix~\ref{app:benchmarks}.



\noindent\textbf{Research questions.}
\textbf{RQ1} (\S\ref{sec:main-results}): How does Adaptive Auto-Harness compare against existing auto-harness systems on open-ended task streams?
\textbf{RQ2} (\S\ref{sec:barriers}): How does Adaptive Auto-Harness address benchmark-specific bottlenecks?
\textbf{RQ3} (\S\ref{sec:ablation}): Does stateful multi-agent evolution with evaluation feedback provide additional gains?
\textbf{RQ4} (\S\ref{sec:navloss}): Can solve-time routing effectively leverage specialized harness branches?
\textbf{RQ5} (\S\ref{sec:hitl}): Can human steering help under insufficient history signal?

\subsection{Comparison with Baselines}
\label{sec:main-results}



\noindent\textbf{Prior systems specialize on one metric cluster.} A-Evolve leads the two pass-rate streams ($45.2\%$ CTF-Dojo, $47.5\%$ FutureX) but covers only $21.1\%$ of PolyBench markets. Meta-Harness leads all three PolyBench metrics ($55.3\%$ Coverage, $50.8\%$ Accuracy, $+320\%$ Return) but falls below the no-evolution Sonnet baseline on FutureX ($29.4\%$ vs.\ $31.0\%$). Base solvers stay below $32.6\%$ PolyBench Coverage on every model, and OctoTools, the frozen human-designed system, places third on PolyBench Return but does not lead any row. The pass-rate cluster and the portfolio cluster therefore sit in different baselines.

\noindent\textbf{Our three variants jointly lead all metrics.} Without HITL intervention, the Full System combines multi-agent evolution with solve-time routing and reaches $97.9\%$ PolyBench Coverage, $80.9\%$ Accuracy, and $50.2\%$ CTF-Dojo Pass. The Multi-agent variant leads FutureX at $49.5\%$, where evolving the right source and tooling matters more than per-task routing. The Adaptive variant leads PolyBench Return at $+352\%$, where matching each market to a specialized strategy matters.

\subsection{Benchmark Bottlenecks}
\label{sec:barriers}

\begin{figure}[!t]
\centering
\includegraphics[width=\linewidth]{figures/l_evo_evidence.pdf}
\vskip -1em
\caption{$\Levo$ evidence across benchmark-specific bottlenecks. PolyBench stresses confidence calibration, FutureX stresses web-retrieval access, and CTF-Dojo stresses payload handling of different file sizes.}
\label{fig:l_evo_evidence}
\vskip -0.5em
\end{figure}

\begin{figure}[!t]
\centering
\includegraphics[width=\linewidth]{figures/l_adapt_evidence.pdf}
\vskip -1em
\caption{$\Ladapt$ evidence across task categories. Curves show cumulative adaptation lift over the baselines, while shaded bands show performance spread across task categories over cycles.}
\label{fig:l_adapt_evidence}
\end{figure}

RQ2 asks whether Adaptive Auto-Harness targets the bottlenecks that limit each benchmark. We organize the analysis around the two losses introduced in \S\ref{sec:framework}: the \emph{evolution loss} $\Levo$, the gap from capabilities the evolver class cannot construct; and the \emph{adaptation loss} $\Ladapt$, the gap from committing to a single harness across heterogeneous tasks.

\noindent\textbf{Evolution bottlenecks ($\Levo$) differ across benchmarks.} Figure~\ref{fig:l_evo_evidence} plots a benchmark-specific stress axis chosen to expose the most predictive capability. The PolyBench panel plots mean stated confidence against market consensus, defined as the implied probability of the favored outcome from Polymarket prices; a well-calibrated harness from our multi-agent variant tracks the diagonal, while single-agent variants stay flat and over-state confidence on low-consensus markets, suggesting that the binding capability is consensus-aware confidence calibration rather than raw prediction skill. The FutureX panel plots pass rate against three retrieval tiers, ranging from offline to date-filtered Wikipedia plus DuckDuckGo to unrestricted DuckDuckGo web search. Pass rate increases monotonically from $34.0\%$ to $47.6\%$ to $57.1\%$, identifying source acquisition rather than reasoning as the binding capability. The CTF-Dojo panel plots pass rate against the largest challenge file size, binned into five tiers from no-payload to $>$1MB. The best single-agent variant declines from $81.8\%$ to $30.4\%$, and the multi-agent variant declines from $90.9\%$ to $39.1\%$ while retaining roughly a 9-point margin throughout, indicating that payload-handling infrastructure becomes the binding capability as inputs grow and that the multi-agent evolver mitigates but does not eliminate this bottleneck. In summary, the binding capability therefore differs by benchmark.

\noindent\textbf{Adaptation bottlenecks ($\Ladapt$) remain after evolution.} Figure~\ref{fig:l_adapt_evidence} reports per-task adaptation lift over evolver cycles, where lift on each task is the adaptative harness score minus the median score across single-committed-harness baselines. The solid curve is the within-cycle mean lift and the shaded band is performance variation across task categories; if a single committed dense harness were sufficient, the mean lift would approach zero as evolution progresses. We observe instead that the mean lift and variation remain positive across all cycles on all three benchmarks. 
% Together with the evolution evidence above, this supports a positive answer to RQ2: each loss term dominates on at least one benchmark, and Adaptive Auto-Harness addresses both. Multi-agent evolution constructs the benchmark-specific capability, and solve-time routing adapts to the appropriate branch for each task.

\noindent\underline{\textbf{Takeaway.}} Although streams expose different bottlenecks, they can be addressed under the same principle: reduce evolution loss for missing capabilities and adaptation loss for task-specific fit.

\subsection{Stateful Multi-Agent Evolution}
\label{sec:ablation}


\begin{figure}[!t]
\centering
\includegraphics[width=\linewidth]{figures/figure_rq2.pdf}
\vskip -0.5em
\caption{Ablations for multi-agent evolution. Removing evaluation feedback or cross-cycle memory degrades performance relative to the full stateful system.}
\label{fig:rq2_stats}
\vskip -0.5em
\end{figure}

\begin{figure}[!t]
\centering
\includegraphics[width=0.9\linewidth]{figures/figure_rq3.pdf}
\vskip -0.5em
\caption{Designed analysis of solve-time routing on an evolved harness tree. We seed one branch per task category, evolve the tree over the stream, and replay every task through every branch. Oracle is the best branch per task, Adapt a category-based routing policy, Naive the \texttt{main} workspace only, and Worst the worst branch per task.}
\label{fig:rq3_stats}
\vskip -0.5em
\end{figure}

\begin{figure*}[!t]
\centering
\begin{minipage}[t]{0.64\textwidth}
\centering
\includegraphics[width=\linewidth]{figures/rq2_demo.pdf}\\[-0.2em]
{\scriptsize (a) Example of the four-phase evolution cycle.}
\end{minipage}\hfill
\begin{minipage}[t]{0.32\textwidth}
\centering
\includegraphics[width=\linewidth]{figures/rq3_demo.pdf}\\[-0.2em]
{\scriptsize (b) Example of the solve-time routing trace.}
\end{minipage}
\vspace{0.25em}

\includegraphics[width=\textwidth]{figures/rq4_demo_horizontal.pdf}\\[-0.2em]
{\scriptsize (c) Example of the Human-in-the-Loop workflow.}
\vskip -0.5em
\caption{Extracted trajectories of the designed multi-agent evolution, agentic routing, and Human-in-the-Loop.}
\label{fig:mechanism_traces}
\vspace{-0.4em}
\end{figure*}


RQ3 asks whether the four-role evolver in section \S\ref{sec:method-multi} (Analyst $\to$ parallel Researchers $\to$ Builder $\to$ Verifier) improves over a single-agent evolver, and whether its state channels matter. Using 100/60/80 tasks on PolyBench/CTF-Dojo/FutureX, Figure~\ref{fig:rq2_stats} compares the full system with variants that remove temporal-reveal feedback or cross-cycle memory. The full system improves over the single-agent evolver on all three benchmarks: $20.3 \to 44.3$ CWR on PolyBench, $38\% \to 43\%$ on CTF-Dojo, and $38\% \to 44\%$ on FutureX. Removing memory causes the broadest degradation, while removing feedback mainly hurts PolyBench, where outcomes resolve after trading. 

\noindent\underline{\textbf{Takeaway.}} The four-role evolver is strongest when paired with persistent state: memory preserves cross-cycle search, and feedback turns resolved outcomes into later evolution signal.



\subsection{Solve-Time Routing on the Harness Tree}
\label{sec:navloss}

RQ4 asks how much adaptation headroom a specialized harness tree exposes, and how much solve-time routing captures. We measure this with a designed analysis rather than end-to-end deployment, isolating the headroom from router quality. On 80/40/58 tasks for PolyBench/CTF-Dojo/FutureX, we seed one branch per task category, evolve the tree over the stream, and replay every task through every branch. We then compare \emph{Oracle} (best branch per task), \emph{Adapt} (category-based routing), \emph{Naive} (\texttt{main} only), and \emph{Worst} (worst per task). As shown in Figure~\ref{fig:rq3_stats}, Adapt improves on both CTF-Dojo and PolyBench, while leaving headroom toward Oracle branch selection. On FutureX, \texttt{main} outperforms Adapt, as web source retrieval ability dominates.
% The answer to RQ4 is therefore positive when branch differences are sufficiently stable to identify, and breaks down on FutureX where rapid source shifts make branch quality non-stationary; the human-steering hooks in \S\ref{sec:hitl} target this case directly.

\noindent\underline{\textbf{Takeaway.}} Harness specialization opens real adaptation headroom, but realized routing captures only part of it, so turning headroom into gain is a separate challenge from constructing the branches.

\subsection{Human Steering for Auto-Harnessing}
\label{sec:hitl}

\begin{figure}[!t]
\centering
\includegraphics[width=\linewidth]{figures/figure_rq4_timeline.pdf}
\caption{FutureX pass-rate lift over four task slices under two human-steering hooks. The orange triangle marks research-phase steering and the purple triangle marks task-board steering. Lift is measured against the no-HITL run.}
\label{fig:rq4_hitl_passrate}
\vskip -0.5em
\end{figure}

RQ5 asks whether human steering helps when history lacks the source or access signal needed for evolution. Since all previous experiment does not apply HITL to ensure fairness, we evaluate this on 100 FutureX tasks and restrict human input to two hooks: \emph{research-phase steering} supplies credentials when research is blocked, and \emph{task-board steering} adds source directions the evolver cannot infer (Figure~\ref{fig:mechanism_traces}c). Figure~\ref{fig:rq4_hitl_passrate} shows the slice-level effect: lift is $0$ on broad polymarket questions, rises to $+5$ on broad search-dependent questions, peaks at $+20$ on the directly targeted finance\&tech slice, and remains $+15$ on adjacent Western-specialty questions. The pattern indicates that HITL helps when it injects the missing external signal, not generic human advice.
% The answer to RQ5 is therefore positive only when the steering targets the specific missing signal, supporting human steering as a complement rather than a substitute for the autonomous evolver.

\noindent\underline{\textbf{Takeaway.}} Human steering is most useful when the missing ingredient is external source knowledge rather than additional autonomous evolution.
